\bookchapter{Κεφάλαιο 40: Μηχανική απαιτήσεων — από την ιδέα σε οδηγίες που η AI $\mu$πορεί να εκτελέσει}

\chapterintro{Αν η $\mu$ηχανική συ$\mu$φραζο$\mu$ένων λύνει το «να καταλαβαίνει η AI το πλαίσιο», η $\mu$ηχανική απαιτήσεων λύνει το «να ξέρει η AI τι ακριβώς πρέπει να κάνει αυτή τη φορά». Εδώ φαίνεται περισσότερο από παντού η διαφορά $\mu$εταξύ ανθρώπων στον προγρα$\mu$$\mu$ατισ$\mu$ό $\mu$ε AI — $\mu$ε τα ίδια εργαλεία και το ίδιο $\mu$οντέλο, άλλοι φτιάχνουν εντυπωσιακά projects και άλλοι $\mu$όνο $\mu$έτρια φτιαγ$\mu$ένα «παιχνιδάκια». Η διαφορά συνήθως δεν είναι στις τεχνικές ικανότητες, αλλά στην ικανότητα να εκφράζει κανείς τις απαιτήσεις του.}

\booksection{Από τη θολή ιδέα στη δο$\mu$η$\mu$ένη απαίτηση}

Κάθε project ξεκινά από $\mu$ια θολή ιδέα — «θέλω να φτιάξω $\mu$ια εφαρ$\mu$ογή καταγραφής εσόδων-εξόδων», «θέλω $\mu$ια προσωπική ιστοσελίδα», «θέλω ένα εργαλείο λίστας υποχρεώσεων». Οι ιδέες αυτές δεν έχουν κανένα πρόβλη$\mu$α καθεαυτές, αλλά αν τις πετάξεις έτσι κι αλλιώς στην AI, η ποιότητα του αποτελέσ$\mu$ατος εξαρτάται εντελώς από την τύχη. Στόχος της $\mu$ηχανικής απαιτήσεων είναι να $\mu$ετατρέπει αυτή τη θολή ιδέα σε $\mu$ια δο$\mu$η$\mu$ένη απαίτηση που η AI $\mu$πορεί να εκτελέσει $\mu$ε ακρίβεια.

Τα καλά νέα: για τη δο$\mu$η$\mu$ένη απαίτηση δεν χρειάζεται να γράψεις PRD δεκάδων σελίδων. Αρκεί να απαντήσεις πέντε ερωτήσεις, που ονο$\mu$άζου$\mu$ε «τα πέντε στοιχεία της απαίτησης»:

\begin{center}
\begin{tabularx}{\linewidth}{|X|X|X|}
\hline
\textbf{Στοιχείο} & \textbf{Ερώτηση} & \textbf{Παράδειγ$\mu$α: εφαρ$\mu$ογή καταγραφής εσόδων-εξόδων} \tabularnewline
\hline
Who & Ποιος θα το χρησι$\mu$οποιεί; & Προσωπική χρήση, καταγραφή καθη$\mu$ερινών εσόδων-εξόδων \tabularnewline
What & Τι θα κάνει; & Μια εφαρ$\mu$ογή που καταγράφει κάθε έσοδο-έξοδο, κάνει $\mu$ηνιαίες συνοψίες και τα απεικονίζει γραφικά \tabularnewline
Why & Γιατί φτιάχνεται; & Θέλω να κατανοήσω τη δο$\mu$ή των εξόδων $\mu$ου και να εντοπίσω περιττές δαπάνες \tabularnewline
How & Χονδρικά πώς φτιάχνεται; & Το frontend σε React, τα δεδο$\mu$ένα τοπικά στον browser, προς το παρόν χωρίς backend και σύνδεση \tabularnewline
Constraint & Ποιοι είναι οι περιορισ$\mu$οί; & Ανάπτυξη από ένα άτο$\mu$ο, ολοκλήρωση $\mu$έσα σε 2 εβδο$\mu$άδες, χωρίς ανάγκη δη$\mu$οσίευσης online \tabularnewline
\hline
\end{tabularx}
\end{center}

\begin{center}\small\color{gray}Τα πέντε στοιχεία της απαίτησης $\mu$ε παραδείγ$\mu$ατα\end{center}

Τα πέντε στοιχεία δεν χρειάζεται να γραφτούν σε επίση$\mu$ο έγγραφο — αρκεί ακό$\mu$α και να τα περάσεις γρήγορα από το $\mu$υαλό σου. Το ζητού$\mu$ενο είναι να αποκτήσεις τη συνήθεια «πρώτα το σκέφτο$\mu$αι καθαρά, $\mu$ετά $\mu$ιλάω». Θα δεις ότι όταν $\mu$πορείς να απαντήσεις καθαρά αυτές τις πέντε ερωτήσεις, η περιγραφή των απαιτήσεών σου γίνεται από $\mu$όνη της δο$\mu$η$\mu$ένη.

\booksection{Η τέχνη της στρω$\mu$ατω$\mu$ένης ανάλυσης των απαιτήσεων}

Για $\mu$εγαλύτερα projects, $\mu$ία περιγραφή απαίτησης δεν αρκεί. Χρειάζεται να σπάσεις τη $\mu$εγάλη απαίτηση σε $\mu$ικρές, ανεξάρτητα επαληθεύσι$\mu$ες υπο-απαιτήσεις. Η διαδικασία αυτή λέγεται «ανάλυση απαιτήσεων», και η βασική της αρχή είναι: κάθε υπο-απαίτηση πρέπει να $\mu$πορεί να ολοκληρωθεί, να επαληθευτεί και να παραδοθεί ανεξάρτητα.

Παράδειγ$\mu$α: για το «φτιάχνω $\mu$ια προσωπική εφαρ$\mu$ογή καταγραφής εσόδων-εξόδων», ένας λογικός τρόπος ανάλυσης είναι:

\begin{enumerate}
  \item Υλοποίηση της βασικής λειτουργίας εγγραφών εσόδων-εξόδων (να προσθέτει, να βλέπει, να διαγράφει εγγραφές)·
  \item Προσθήκη λειτουργίας κατηγοριών (ετικέτες όπως εστίαση, $\mu$ετακινήσεις σε κάθε εγγραφή)·
  \item Υλοποίηση $\mu$ηνιαίας σύνοψης (στατιστικά εσόδων-εξόδων ανά $\mu$ήνα)·
  \item Προσθήκη οπτικοποίησης δεδο$\mu$ένων (απεικόνιση της δο$\mu$ής εξόδων $\mu$ε διαγρά$\mu$$\mu$ατα).
\end{enumerate}

Κάθε βή$\mu$α είναι $\mu$ια λειτουργική βελτίωση που $\mu$πορεί να τρέξει και να επαληθευτεί ανεξάρτητα.

Ο χρυσός κανόνας της ανάλυσης απαιτήσεων: στο τέλος κάθε βή$\mu$ατος, το έργο πρέπει να είναι "κάτι που δουλεύει"—όσο τραχύ κι αν είναι. Ποτέ $\mu$ην αφήνεις το έργο σε κατάσταση "όλα στη $\mu$έση, τίποτα δουλεύει".

Αυτός ο τρόπος ανάλυσης έχει δύο οφέλη. Πρώτον, σου επιτρέπει να βλέπεις διαρκώς πρόοδο — $\mu$ε κάθε ολοκληρω$\mu$ένο βή$\mu$α η εφαρ$\mu$ογή σου αποκτά $\mu$ια λειτουργία ακό$\mu$α, και αυτή η θετική ανατροφοδότηση είναι ζωτικής ση$\mu$ασίας για projects $\mu$εγάλης διάρκειας (όλοι χρειαζό$\mu$αστε λίγη αίσθηση επιτυχίας). Δεύτερον, κρατά κάθε εργασία της AI αρκετά $\mu$ικρή, $\mu$ειώνοντας τον κίνδυνο «παρανόησης από την AI» — όσο $\mu$ικρότερη η εργασία, τόσο ακριβέστερη η κατανόηση και τόσο υψηλότερη η ποιότητα του αποτελέσ$\mu$ατος.

\booksection{Πρότυπο και πρακτική συγγραφής PRD}

Για $\mu$εγαλύτερα projects (π.χ. $\mu$ια full-stack εφαρ$\mu$ογή $\mu$ε frontend και backend και πολλαπλά λειτουργικά modules), τα πέντε στοιχεία $\mu$πορεί να $\mu$ην φτάνουν. Τότε χρειάζεσαι ένα πιο επίση$\mu$ο έγγραφο απαιτήσεων project (PRD). Ένα ση$\mu$είο που παρερ$\mu$ηνεύεται εύκολα: το PRD δεν γράφεται για να το διαβάσει η AI — γράφεται για σένα, για να τακτοποιήσεις τη σκέψη σου· $\mu$όνο τότε $\mu$πορείς να δώσεις στην AI πιο ακριβείς οδηγίες.

Ένα καλό PRD δεν χρειάζεται να είναι $\mu$ακροσκελές, αλλά πρέπει να περιλα$\mu$βάνει έξι $\mu$έρη: σύνοψη project (σε $\mu$ία παράγραφο τι ακριβώς κάνει), λίστα λειτουργιών (τι είναι κάθε ση$\mu$είο-λειτουργία), τεχνική προσέγγιση (ποιο τεχνολογικό stack, χονδρική αρχιτεκτονική), $\mu$οντέλο δεδο$\mu$ένων (οι βασικές δο$\mu$ές δεδο$\mu$ένων), περιγραφή διεπαφής (πώς $\mu$οιάζουν οι κύριες σελίδες), πλάνο ανάπτυξης (σε πόσα βή$\mu$ατα ολοκληρώνεται).

\begin{lstlisting}
# PRD εφαρ$\mu$ογής προσωπικού λογιστηρίου
## Επισκόπηση έργου
Μια ελαφριά εφαρ$\mu$ογή προσωπικού λογιστηρίου, $\mu$ε καταγραφή εσόδων/εξόδων, κατηγορίες, $\mu$ηνιαία σύνοψη και οπτικοποίηση δεδο$\mu$ένων.
## Λίστα λειτουργιών
1. Καταγραφή εσόδων/εξόδων: ποσό, τύπος (έσοδο/έξοδο), κατηγορία, ση$\mu$είωση, η$\mu$ερο$\mu$ηνία
2. Διαχείριση κατηγοριών: έτοι$\mu$ες οι φαγητό, $\mu$ετακινήσεις, ψώνια, διασκέδαση, άλλα
3. Μηνιαία σύνοψη: συνολικά έσοδα, έξοδα και υπόλοιπο ανά $\mu$ήνα
4. Οπτικοποίηση εξόδων: πίτα $\mu$ε τα ποσοστά κάθε κατηγορίας
5. Εξαγωγή δεδο$\mu$ένων: σε CSV
## Τεχνικός σχεδιασ$\mu$ός
- Frontend: React + TypeScript + Tailwind CSS
- Αποθήκευση: localStorage (χωρίς backend)
- Γραφή$\mu$ατα: $\mu$ε Chart.js
- Build: Vite
## Μοντέλο δεδο$\mu$ένων
Record {
  id: string
  amount: number
  type: 'income' | 'expense'
  category: string
  note: string
  date: string  // YYYY-MM-DD
}
## Πλάνο ανάπτυξης
Step 1: βασική καταγραφή (CRUD)
Step 2: διαχείριση κατηγοριών
Step 3: $\mu$ηνιαία σύνοψη
Step 4: οπτικοποίηση δεδο$\mu$ένων
Step 5: εξαγωγή δεδο$\mu$ένων
\end{lstlisting}

Πρόσεξε το στυλ συγγραφής αυτού του PRD: σύντο$\mu$ες λίστες και δο$\mu$η$\mu$ένη $\mu$ορφή, όχι φουσκω$\mu$ένα κεί$\mu$ενα. Κάθε ση$\mu$είο-λειτουργία έχει ξεκάθαρο ορισ$\mu$ό, και το $\mu$οντέλο δεδο$\mu$ένων περιγράφεται $\mu$ε απλούς τύπους. Ένα PRD σε αυτό το στυλ όχι $\mu$όνο το διαβάζεις εσύ καθαρά, αλλά όταν το δώσεις στην AI, εκείνη κατανοεί $\mu$ε ακρίβεια κάθε ση$\mu$είο της απαίτησης.

\booksection{Στην πράξη: πλήρης πορεία από την ιδέα στο σχέδιο}

Δεν έχει νόη$\mu$α να $\mu$ένου$\mu$ε στα θεωρητικά: ας πάρου$\mu$ε ένα πλήρες παράδειγ$\mu$α και ας διανύσου$\mu$ε όλη τη διαδικασία της $\mu$ηχανικής απαιτήσεων από την αρχή ως το τέλος. Ας πού$\mu$ε ότι η ιδέα σου είναι: «θέλω να φτιάξω έναν καταγραφέα $\mu$άθησης που να παρακολουθεί τι έ$\mu$αθα κάθε $\mu$έρα».

Πρώτο βή$\mu$α: περνάς τα πέντε στοιχεία της απαίτησης.

\begin{itemize}
  \item Who: προσωπική χρήση.
  \item What: ένα εργαλείο που καταγράφει τι $\mu$ελέτησες κάθε $\mu$έρα, παρακολουθεί τον χρόνο $\mu$ελέτης και επιτρέπει ανασκόπηση της προόδου.
  \item Why: θέλω να αποκτήσω συνήθεια συνεχούς $\mu$άθησης και να βλέπω τι έχει συσσωρευτεί.
  \item How: πρώτα $\mu$ια web έκδοση, $\mu$ε τα δεδο$\mu$ένα τοπικά στον browser.
  \item Constraint: ανάπτυξη από το $\mu$ηδέν, ολοκλήρωση των βασικών λειτουργιών $\mu$έσα σε 1 εβδο$\mu$άδα.
\end{itemize}

Δεύτερο βή$\mu$α: αναλύεις τις απαιτήσεις.

\begin{enumerate}
  \item Βή$\mu$α 1: υλοποίηση της λειτουργίας «καταγραφή εγγραφών $\mu$άθησης» (η$\mu$ερο$\mu$ηνία, περιεχό$\mu$ενο, διάρκεια)·
  \item Βή$\mu$α 2: προσθήκη «στατιστικών χρόνου $\mu$ελέτης» (σή$\mu$ερα, εβδο$\mu$αδιαίο σύνολο)·
  \item Βή$\mu$α 3: υλοποίηση προβολής «η$\mu$ερολογίου $\mu$άθησης» (όπως το γράφη$\mu$α contributions του GitHub)·
  \item Βή$\mu$α 4: προσθήκη λειτουργίας «ανασκόπησης $\mu$άθησης» (τυχαία ε$\mu$φάνιση εγγραφών από κάποια προηγού$\mu$ενη $\mu$έρα, για επανάληψη).
\end{enumerate}

Τρίτο βή$\mu$α: γράφεις ένα σύντο$\mu$ο PRD. Δεν χρειάζεται να είναι $\mu$ακρύ — η ανάλυση παραπάνω αποτελεί ήδη το κύριο σώ$\mu$α του PRD. Το τακτοποιείς σε ένα αρχείο .md, το βάζεις στον ριζικό φάκελο του project, και λειτουργεί ως κοινό ση$\mu$είο αναφοράς για σένα και την AI.

Τέταρτο βή$\mu$α: ξεκινάς την ανάπτυξη του πρώτου βή$\mu$ατος. Η οδηγία προς την AI $\mu$πορεί να είναι:

\begin{lstlisting}
Χρειάζο$\mu$αι την πρώτη λειτουργία ενός ανιχνευτή $\mu$άθησης—καταγραφή $\mu$αθη$\mu$άτων.
Ο χρήστης $\mu$πορεί να προσθέσει $\mu$ια εγγραφή $\mu$ε η$\mu$ερο$\mu$ηνία, περιεχό$\mu$ενο και διάρκεια (λεπτά).
Τα δεδο$\mu$ένα αποθηκεύονται σε localStorage.
Frontend $\mu$ε React + TypeScript.
Δώσε πλήρη υλοποίηση.
\end{lstlisting}

Πρόσεξε τα χαρακτηριστικά αυτής της οδηγίας: περιγράφει $\mu$όνο τον στόχο του τρέχοντος βή$\mu$ατος (τη λειτουργία καταγραφής), χωρίς να $\mu$πλέκει άλλες λειτουργίες· ορίζει ξεκάθαρα τους τεχνικούς περιορισ$\mu$ούς (React, TypeScript, localStorage)· δίνει αρκετές πληροφορίες ώστε η AI να ξέρει «τι να κάνει» και «πώς να το κάνει», ενώ τη συγκεκρι$\mu$ένη υλοποίηση την αφήνει στην AI να αποφασίσει. Αυτή είναι $\mu$ια υψηλής ποιότητας οδηγία προγρα$\mu$$\mu$ατισ$\mu$ού $\mu$ε AI, κατεργασ$\mu$ένη $\mu$έσω $\mu$ηχανικής απαιτήσεων.

\begin{tcbpractice}{Πά$\mu$ε στα πρακτικά: φτιάξε $\mu$ια κάρτα πέντε στοιχείων για το project που σε «καίει»}
Διάλεξε ένα $\mu$ικρό project που το ζυ$\mu$ώνεις καιρό στη σκέψη σου και πέρασε 5 λεπτά γε$\mu$ίζοντας $\mu$ια κάρτα: Who / What / Why / How / Constraint, από $\mu$ία πρόταση το καθένα. Μετά χώρισέ το σε 3-5 βή$\mu$ατα, ώστε «κάθε βή$\mu$α να είναι χρήσι$\mu$ο από $\mu$όνο του». Θα δεις: $\mu$όλις γράψεις την ιδέα, από «ονειροπόλη$\mu$α» γίνεται «λίστα υποχρεώσεων».

\end{tcbpractice}

\begin{tcbsummary}{Σύνοψη κεφαλαίου}
\begin{itemize}
  \item Τα πέντε στοιχεία της απαίτησης: Who / What / Why / How / Constraint — σκέψου καθαρά πριν $\mu$ιλήσεις
  \item Χρυσός κανόνας ανάλυσης απαιτήσεων: $\mu$ετά από κάθε βή$\mu$α το project είναι «κάτι που δουλεύει»
  \item Το PRD γράφεται για σένα: σύνοψη, λίστα λειτουργιών, τεχνική προσέγγιση, $\mu$οντέλο δεδο$\mu$ένων, διεπαφή, πλάνο
  \item Καλή οδηγία = περιγράφεις $\mu$όνο το τρέχον βή$\mu$α + ξεκάθαροι περιορισ$\mu$οί + λες «τι να κάνει», όχι «πώς να το κάνει»
  \item Πρότυπο πλήρους πορείας: πέντε στοιχεία ανάλυση PRD οδηγία, σε τέσσερα βή$\mu$ατα
\end{itemize}

\end{tcbsummary}
